\textbf{Infográfico: How a Bottle Is Recycled}

\begin{quote}
\textbf{Good job!} You tossed that plastic bottle into the recycling bin rather than straight into the trash. But do you know what will happen to that bottle now?

\textbf{1. Collection.} Bottles and other recyclable materials are picked up from homes, schools, and businesses and taken to a recycling facility.

\textbf{2. Sorting.} Bottles and containers are sorted based on what they're made of and sometimes what color they are. The items then go through a machine that shreds them.

\textbf{3. Melting.} The shreds are washed to remove any impurities, pieces of labels, or leftover contents. They are then dried and melted.

\textbf{4. Flaking.} The melted plastic is processed into flakes, which can be made into a wide variety of plastic products.

\textbf{5. Next applications.} The flakes can be spun into a very fine, threadlike material in a process similar to how cotton is made. Companies use those ``threads'' to make a wide variety of items, including carpets, clothing, or filling for jackets and quilts. Depending on the type of plastic, they can also be made into new (well, sort of new!) plastic bottles.

\textbf{Amount of trash the average person produces in a single day:} 4.43 pounds.

\textbf{Percentage of trash that's recycled:} 34.1\%.
\end{quote}

\textit{Sources: U.S. EPA; Earth911.}

\textbf{COLUNI 2018 -- Questão 22}\par\medskip Em relação às informações apresentadas no texto, é CORRETO afirmar:\begin{enumerate}
\item Na etapa 1, os consumidores levam as garrafas vazias até o local de reciclagem.
\item Na etapa 2, as garrafas são selecionadas por cor e tipo de material e depois lavadas.
\item Na etapa 3, os resíduos são retirados das garrafas e o plástico é derretido.
\item Na etapa 5, a produção dos flocos é semelhante à produção do algodão doce.

CIÊNCIAS – QUESTÕES DE 23 A 30\end{enumerate}
